\section{Проводная сеть: TWAI/CAN и RS-485}\label{les:wired-buses}
\lessonintro
  {разделять контроллер и физический интерфейс и связать два узла}
  {уроки 2 и 15}
  {сообщения между двумя ESP32-S3 по кабелю}
  {две платы и два CAN- либо RS-485-трансивера с логикой 3,3 В}

\subsection{Контроллер, трансивер и протокол}
Дифференциальный приёмник сравнивает две линии: одинаково наведённая помеха меньше влияет
на результат. Это не отменяет ограничений по длине, скорости, земле и согласованию.
При разных источниках питания проверьте допустимую разность потенциалов; удалённым узлам
могут понадобиться развязка и защита, которых нет на простой макетке.

\begin{center}\small
\begin{tabularx}{\textwidth}{@{}l X X@{}}\toprule
 & TWAI / Classical CAN & UART через RS-485\\\midrule
В чипе & контроллер кадров и арбитража & UART, направление задаёт программа\\
Снаружи & CAN-трансивер & RS-485-трансивер\\
Данные & ID и до 8 байтов & поток байтов, границы задаёт приложение\\
\bottomrule\end{tabularx}\end{center}
TWAI ESP32-S3 поддерживает Classical CAN, а не CAN FD. RS-485 --- электрический интерфейс,
он сам не добавляет адреса, контрольную сумму и команды Modbus.

\subsection{Два узла CAN}
Используйте отдельную учебную шину. GPIO16 идёт на TXD трансивера, RXD на GPIO17.
Соедините CANH с CANH, CANL с CANL, сигнальные земли в пределах стенда объедините.
На двух физических концах ставят по 120 Ом между CANH и CANL, а не на каждом узле.
У модулей резисторы могут быть уже установлены. Проверьте питание модели и допустимые уровни RXD.

\codecap{Кадры Classical CAN со счётчиком}
\codeinput{l32_twai}

На второй плате измените ID на \code{0x102}. Скорость обоих узлов 500 кбит/с.
Успех \code{twai\_transmit} означает постановку в очередь, а не доставку приложению.
Один узел в нормальном режиме не получит ACK от самого себя. Подтвердите приём другой платой.
Для гарантии обработки команды добавьте ответ приложения с номером сообщения.
При BUS OFF сначала ищут неисправность линии, затем выполняют восстановление по драйверу.

\subsection{Два узла RS-485}
У полудуплексного модуля соедините DE и активный LOW вход /RE с GPIO4. DI к GPIO16,
RO к GPIO17. HIGH включает передатчик, LOW --- приёмник. Линии соединяйте по документации:
обозначения A/B у производителей могут различаться. Согласование ставят на концах;
смещение задают в одном рассчитанном месте либо используют трансивер с нужным failsafe.

\codecap{Учебные запрос PING и ответ PONG}
\codeinput{l32_rs485}

На второй плате \code{MASTER=false}. Только ведущий начинает обмен, второй отвечает после запроса.
\code{flush()} ждёт передачи последнего стоп-бита до выключения DE.
Это учебный текстовый протокол, не Modbus RTU. В реальной сети нужны адреса, контроль целостности,
тайм-аут, ограниченный повтор и правило предотвращения коллизий. На длинной линии важны
топология и длина ответвлений: набор соединённых перемычек не моделирует любой кабель.

\begin{trouble}
CAN без ACK: нет второго рабочего узла, неверна скорость или проводка.
RS-485 молчит: неверна полярность пары, DE остался включён либо /RE не активен.
Проверяйте физическую линию до разбора содержимого кадров.
\end{trouble}
\begin{summary}
Контроллер формирует данные, трансивер работает с кабелем. CAN имеет арбитраж и кадры,
а поверх RS-485 протокол выбирают отдельно. Очередь и доставка --- разные события.
\end{summary}

\begin{selfcheck}
\begin{enumerate}
\item Почему нельзя соединить GPIO с CANH напрямую?
\item Где ставят согласующие резисторы?
\item Почему успешный \code{twai\_transmit} не доказывает доставку?
\end{enumerate}
\end{selfcheck}
\needspace{20\baselineskip}
\begin{tasks}
\begin{enumerate}
\item \lv{1} Найдите контроллер, трансивер и кабель на схеме.
\item \lv{1} Настройте два CAN-узла с разными ID и подтвердите приём.
\item \lv{2} Передайте 16-битную температуру с масштабом 0,01 градуса и опишите порядок байтов.
\item \lv{2} Повторите PING/PONG через RS-485 и измерьте время ответа.
\item \lv{3} Добавьте номер запроса, контрольную сумму и ограниченные повторы; проверьте отключение второго узла.
\item \lv{4} Спроектируйте сеть трёх датчиков: адреса, опрос, согласование, тайм-ауты и реакции на неисправности.
\end{enumerate}
\end{tasks}

\paragraph{Документация для продолжения}\href{https://docs.espressif.com/projects/esp-idf/en/v5.4/esp32s3/api-reference/peripherals/twai.html}{Официальное руководство по теме}.
\href{https://docs.espressif.com/projects/esp-idf/en/v5.4/esp32s3/api-reference/peripherals/uart.html}{UART и RS-485}.
